\uuid{IdfN}
\exo7id{4794}
\titre{exo7 4794}
\auteur{quercia}
\organisation{exo7}
\datecreate{2010-03-16}
\isIndication{false}
\isCorrection{false}
\chapitre{Topologie}
\sousChapitre{Topologie des espaces vectoriels normés}
\cpge{mpsi-mp2i.fonctions-deux-variables.a, pcsi.fonctions-deux-variables.a, ptsi.fonctions-deux-variables.a, mp-mpi.topologie-espaces-normes.e, pc.espaces-vectoriels-normes.e, psi.espaces-vectoriels-normes.e, pt.fonctions-deux-trois-variables.b}
\module{Analyse}
\niveau{L2}
\difficulte{}

\contenu{
\question{
Trouver une fonction $f : {\R^2} \to \R$ discontinue en $(0,0)$
mais telle que pour tous $\alpha,\beta \in \R$,
$f(\alpha t,\beta t) \xrightarrow[t\to0]{} 0$.
}
}
